\begin{frame}{Presiones de gasto durante la próxima década}
\small
\begin{itemize}
  \item Se puede lograr un recorte de más de dos puntos del PIB revirtiendo aumentos en inversión, subsidios a los combustibles y otras transferencias.
  \item Sin embargo, persistirán presiones adicionales en salud y pensiones.
\end{itemize}
\centering
\begin{tabular}{lc}
\toprule
\textbf{Rubro} & \textbf{Aumento 2026--2036 (\% del PIB)} \\
\midrule
Salud & 1,8 \\
Colpensiones & 0,6 \\
Asignaciones de retiro (militares y Policía) & 0,5 \\
\midrule
\textbf{Total} & \textbf{2,9} \\
\bottomrule
\end{tabular}

\vfill
\scriptsize Fuente: MFMP 2025. Esta tabla ya cubre un horizonte posterior a 2025.
\end{frame}

\begin{frame}{Contener el gasto en salud y pensiones}
\begin{itemize}
  \item La consolidación fiscal difícilmente será posible sin contener las presiones de gasto en salud y pensiones.
  \item Para \textbf{salud} se requiere:
  \begin{itemize}
    \item buscar fuentes de financiación por fuera del PGN;
    \item retornar a un plan de beneficios explícitos.
  \end{itemize}
  \item Para \textbf{pensiones} se requiere:
  \begin{itemize}
    \item aplicar reformas paramétricas;
    \item gravar las pensiones con el impuesto de renta;
    \item revisar el mandato constitucional de una pensión mínima de un salario mínimo.
  \end{itemize}
\end{itemize}
\end{frame}

\begin{frame}{El traslado de competencias debe acompañar los recursos}
\begin{itemize}
  \item A las presiones de gasto se suma el aumento proyectado del \textbf{SGP}, que podría acercarse a dos puntos del PIB por el Acto Legislativo 03 de 2024.
  \item La presión podría neutralizarse si la Ley de Competencias garantiza correspondencia entre los recursos adicionales y el costo de las funciones trasladadas a las entidades territoriales.
  \item Conviene evaluar si la ley debe tramitarse antes de contar con esa correspondencia y con una senda fiscal financiable.
\end{itemize}
\end{frame}

\begin{frame}{El SGP amplifica el recaudo requerido}
\begin{itemize}
  \item Una reducción de al menos 3 puntos del PIB en el déficit primario requerirá aumentar los ingresos.
  \item La fórmula del SGP implica que cada 100 pesos de ingresos corrientes adicionales generan eventualmente al menos 25 pesos de transferencias.
  \item Por tanto, cerrar 3 puntos del PIB mediante ingresos exige un recaudo bruto cercano a \textbf{4 puntos del PIB}.
  \item Si la Ley de Competencias crea obligaciones adicionales no financiadas, el recaudo requerido sería mayor.
\end{itemize}
\end{frame}

\begin{frame}{Resumiendo}
\begin{itemize}
  \item La \textbf{consolidación fiscal} requiere:
  \begin{itemize}
    \item recortar inversión y otras transferencias;
    \item contener el gasto en pensiones y salud;
    \item revisar el calendario de la Ley de Competencias;
    \item reformular el SGP como porcentaje del PIB;
    \item marchitar beneficios tributarios;
    \item reformar administrativamente la DIAN.
  \end{itemize}
  \item La gradualidad requerirá una \textbf{nueva regla fiscal}.
\end{itemize}
\end{frame}

\begin{frame}{¿Por qué siguen prestando los inversionistas?}
\small
\begin{itemize}
  \item Es una combinación de restricciones institucionales e incentivos:
  \begin{itemize}
    \item \textbf{AFP y bancos domésticos}: no tienen una alternativa local de escala comparable a los TES.
    \item \textbf{Inversionistas extranjeros}: aprovechan el \textit{carry} mientras no haya una señal clara de incumplimiento.
    \item \textbf{Represión financiera}: algunos anticipan tasas reales negativas y una licuación gradual de la deuda, no un incumplimiento abrupto.
  \end{itemize}
  \item Si el Banco de la República no monetiza, el ajuste o la reestructuración ganan relevancia; si cede, se debilita el ancla que sostiene la demanda por TES.
  \item La presión para ajustar puede terminar llegando con un choque externo: menor precio del petróleo, mayores tasas globales o una nueva rebaja soberana.
\end{itemize}
\end{frame}

\begin{frame}{Lecciones de casos comparados: mecanismos}
\small
\centering
\begin{tabular}{l p{3.2cm} p{2.5cm} p{2.4cm}}
\toprule
\textbf{Caso} & \textbf{Mecanismo} & \textbf{Catalizador} & \textbf{Costo} \\
\midrule
Ecuador 2000 & Dolarización forzada e incumplimiento & Crisis bancaria y choques externos & Devastador \\
Ecuador 2020 & Reestructuración y acuerdo con el FMI & COVID-19 & Moderado \\
Costa Rica 2018--24 & Reforma fiscal negociada & Rebaja a alto rendimiento & Alto, pero manejable \\
\textbf{Colombia (?)} & \textbf{Por determinar} & \textbf{Por determinar} & \textbf{Por determinar} \\
\bottomrule
\end{tabular}

\vfill
\scriptsize Fuente: elaboración propia con base en FMI, OCDE y BFI.
\end{frame}

\begin{frame}{Lecciones de casos comparados: implicaciones}
\begin{itemize}
  \item En los casos revisados, el ajuste no surgió sin presión externa. La diferencia fue la capacidad institucional para \textbf{negociarlo} en lugar de \textbf{sufrirlo}.
  \item \textbf{Ecuador}: el ajuste políticamente imposible se volvió inevitable; el incumplimiento de 2008 respondió a falta de voluntad de pago, no a incapacidad fiscal inmediata.
  \item \textbf{Costa Rica}: la reforma de 2018 necesitó como catalizador una rebaja a grado especulativo.
  \item Para Colombia, la pregunta es si las instituciones permitirán anticipar el ajuste antes de que el mercado lo imponga.
\end{itemize}
\end{frame}